\section{Experiments: Training-Light Mechanism}\label{sec:results-tl}

\subsection{Layer-position ablation}\label{sec:bandablation}

Four band-LoRA arms, matched trainable-parameter budget (early 0--15
$r{=}8$; mid 10--25 $r{=}8$, a pre-existing production checkpoint; late
14--29 $r{=}8$; full-depth 0--29 $r{=}4$, rank lowered to stay near the
16-layer arms' 3.93M-parameter budget), 4000 steps, held-out 30-subject
dev set $\times$ 2 seeds ($n{=}60$/arm), same defect-aware judge protocol
as the training-free experiments above. The reused mid checkpoint later
proved \emph{not} protocol-matched to the other arms (different
training-manifest revision; a style prefix in its generation prompts) --
a confound we discovered, quantified, and resolve below. \textbf{Pre-registered decision rule,
locked before the three missing arms were run}: a band counts as
specially favored on identity-match only if its clean-yes point estimate
is highest \emph{and} its 95\% CI does not overlap the worst arm's.

\begin{table}[t]
\centering\small
\begin{tabular}{@{}lrrr@{}}
\toprule
Arm & clean-yes & defect & clean-usable \\
\midrule
early 0--15 ($r{=}8$) & 15.0\% [8,26] & 8.3\% [4,18] & 45.0\% [33,58] \\
mid 10--25 ($r{=}8$)$^\dagger$ & 21.7\% [13,34] & \textbf{3.3\% [1,11]} & 55.0\% [42,67] \\
late 14--29 ($r{=}8$) & 21.7\% [13,34] & 18.3\% [11,30] & 46.7\% [35,59] \\
full-depth ($r{=}4$) & 23.3\% [14,35] & \textbf{28.3\% [19,41]} & 43.3\% [32,56] \\
\bottomrule
\end{tabular}
\caption{Band-position ablation on the first dev set, $n{=}60$/arm,
Wilson 95\% CI. $^\dagger$The mid row was measured under a non-shared
protocol; a protocol-matched re-run yields defect 30.0\% [20,43] (see
text).}
\label{tab:bandablation}
\end{table}

\textbf{Primary metric result: null, replicated.} All four clean-yes
CIs overlap (highest, full-depth, 23.3\% [14,35]; lowest, early, 15.0\%
[8,26]) -- under the pre-registered rule, band position is \emph{not}
specially favored for identity-match. The a priori hypothesis (``10--25
is the identity band'') fails on its own designed test. A full re-run
of all four checkpoints on a second, disjoint held-out 30-subject set
(zero retraining, identical judge) reproduces the null: clean-yes
18.3--25.0\%, all CIs overlapping. Identity capacity is depth-uniform
in every cell we measured.

\textbf{Secondary observation, refuted by our own audit.} On the first
set, defect rates appeared to form a CI-decisive gradient
(3.3\%$\to$8.3\%$\to$18.3\%$\to$28.3\%; subject-clustered mid 6.7\%
[0,17] vs.\ full-depth 43.3\% [27,60], non-overlapping). Three audits
then killed it. \emph{(i) Cross-set replication fails}: on the second
set the interior ranking inverts (late becomes best, 8.3\%; mid
second-worst, 20.0\%) and mid-vs.-full-depth clustered CIs overlap
(36.7\% [20,53] vs.\ 46.7\% [30,63]). \emph{(ii) The mid arm was
double-confounded}: a completed 2$\times$2 factorial (checkpoint
$\times$ prompt-prefix, $n{=}60$/cell) closes the attribution instead
of bounding it -- the prefix effect is checkpoint-specific, not
portable ($+13.3$pp within run\_mix1, $+0$pp within run\_mix2), and the
only clustered-CI-disjoint marginal is the checkpoint swap at
prefix-on ($+26.7$pp); clean-yes stays flat across all four cells
(16.7--21.7\%). \emph{(iii) A confound-free mid} (same recipe,
matched manifest, no prefix) scores 30.0\% [20,43] defect on the
\emph{original} set -- clustered 53.3\% [37,70], disjoint from the
confounded 6.7\% [0,17] -- and is cross-set stable (26.7\% on the
second set). Under a uniform protocol every arm sits in an 8--30\%
defect band with no stable position ordering; full-depth is numerically
worst on both sets but never CI-separated from mid. A seeded random
audit of judge flags (10 of 17 full-depth flags, all mid flags,
negative controls) additionally found only 30--70\% of flags
correspond to real damage depending on severity threshold (zero false
negatives on controls) -- raw judge rates overstate severe-defect
rates. What survives is the null: layer position determines neither
identity capacity nor, as far as our evidence goes, defect rate.

\subsection{Does an identity-supervised objective separate the
defect?}\label{sec:auxloss}

If the collage defect were a training deficiency rather than a
structural entanglement, adding an explicit identity signal to the
objective should buy identity without the defect. We test this directly.
From a config-matched MSE-only baseline (band 10--25, $r{=}8$, mix3
manifest, seed 0, 1500 steps), two arms add one auxiliary loss and
nothing else: a decode-space CCIP identity loss (weight 0.05) and a
perceptual loss (weight 0.1). Evaluated on the full held-out
30-subject set $\times$ 2 seeds ($n{=}60$/arm) with the same
defect-aware judge, neither helps (Table~\ref{tab:auxloss}): both leave
clean-yes at or below the MSE baseline while \emph{raising} the defect
rate -- the identity-loss arm to 26.7\% and the perceptual arm to
40.7\%, against the baseline's 13.3\%. The most direct training-side
repair moves the tradeoff the wrong way rather than dissolving it. This
is the fifth matched-control repair, and the only training-side one, to
trade identity for coherence rather than separate them -- consistent
with a defect entangled with identity, not one caused by weak
supervision.

\begin{table}[t]
\centering\small
\begin{tabular}{@{}lrr@{}}
\toprule
Arm ($n{=}60$) & clean-yes & defect \\
\midrule
MSE baseline & \textbf{16.7\% [9,28]} & \textbf{13.3\% [7,24]} \\
+ CCIP identity loss & 15.0\% [8,26] & 26.7\% [17,39] \\
+ perceptual loss & 10.2\% [5,20] & 40.7\% [29,53] \\
\bottomrule
\end{tabular}
\caption{Auxiliary-loss pilot, held-out $n{=}60$/arm, Wilson 95\% CI. An
explicit identity (CCIP) or perceptual objective does not reduce the
composition defect; both raise it without improving identity match.}
\label{tab:auxloss}
\end{table}
